
\documentclass[oneside]{diretrizes}            % Imprimir apenas frente
%\documentclass[doubleside]{diretrizes}        % Imprimir frente e verso

% Importações de pacotes
\usepackage[alf, abnt-emphasize=bf, recuo=0cm, abnt-etal-cite=2, abnt-etal-list=0, abnt-repeated-author-omit=yes]{abntex2cite}  % Citações padrão ABNT
\usepackage{listings}

\usepackage[utf8]{inputenc}                         % Acentuação direta
\usepackage[T1]{fontenc}                            % Codificação da fonte em 8 bits
\usepackage{graphicx}                               % Inserir figuras
\usepackage{amsfonts, amssymb, amsmath}             % Fonte e símbolos matemáticos
\usepackage{booktabs}                               % Comandos para tabelas
\usepackage{verbatim}                               % Texto é interpretado como escrito no documento
\usepackage{multirow, array}                        % Múltiplas linhas e colunas em tabelas
\usepackage{indentfirst}                            % Endenta o primeiro parágrafo de cada seção.
\usepackage{microtype}                              % Para melhorias de justificação?
\usepackage[algoruled, portuguese]{algorithm2e}     % Escrever algoritmos
\usepackage{float}                                  % Utilizado para criação de floats
\usepackage{times}                                  % Usa a fonte Times
\usepackage{color}
\usepackage{xcolor}
\usepackage{listings}

% Definindo novas cores
\definecolor{verde}{rgb}{0.25,0.5,0.35}
\definecolor{jpurple}{rgb}{0.5,0,0.35}
% Configurando layout para mostrar codigos Java
\lstset{
  language=Java,
  basicstyle=\ttfamily\small, 
  keywordstyle=\color{jpurple}\bfseries,
  stringstyle=\color{red},
  commentstyle=\color{verde},
  morecomment=[s][\color{blue}]{/**}{*/},
  extendedchars=true, 
  showspaces=false, 
  showstringspaces=false, 
  numbers=left,
  numberstyle=\tiny,
  breaklines=true, 
  backgroundcolor=\color{cyan!10}, 
  breakautoindent=true, 
  captionpos=b,
  xleftmargin=0pt,
  tabsize=4
}

% Inclui o preâmbulo do documento
%
% Documento: Preâmbulo
%

\usepackage{helvet}
\renewcommand{\familydefault}{\sfdefault}

\instituicao{Instituto Federal de Educação, Ciência e Tecnologia da Paraíba}
\abreviatura{IFPB}
\departamento{Campus Cajazeiras}
\localapresentacao{Cajazeiras}
\programa{Curso Superior de Tecnologia}
\modalidade{em Análise e Desenvolvimento de Sistemas}
\nomeautor{Nome do Autor}
\titulotb{Título do trabalho}
%\subtitulo{Aprovado pelo Colegiado em XX/XX/XXXX}
\anoapresentacao{2019}
\grau{Tecnólogo em Análise e Desenvolvimento de Sistemas}
\dataapresentacao{2023}
\mesapresentacao{data da apresentação}

%Dados Orientador
\orientador{Francisco Paulo de Freitas Neto}
%\coorientador{Nome do Coorientador}
\instOrientador{IPB}
\departamentoorientador{Campus Cajazeiras}
\titulacaoorientador{Prof. Me.}

%Dados Coorientador
\coorientador{Nome do Coorientador, se houver}
\instCoorientador{IFPB}
\departamentocoorientador{Campus Cajazeiras}
\titulacaocoorientador{Prof. Me.}

%Dados Examinador 1
\nmexamum{Nome do examinador 1}
\instexamum{Instituição do avaliador}
\departamentoexamum{Campus do avaliador}
\titulacaoexamum{Prof. Me.}

%Dados Examinador 2
\nomeexamdois{Nome do examinador 2}
\instexamdois{Instituição do avaliador}
\departamentoexamdois{Campus do avaliador}
\titulacaoexamdois{Prof. Me.}

\usepackage[a4paper, top=3.0cm, bottom=2.0cm, left=3.0cm, right=2.0cm]{geometry} %teste margem
%\usepackage{titlesec} %teste margem

\newcommand{\blue}[1]{\textcolor{blue}{#1}} % proposed
\newcommand{\red}[1]{\textcolor{red}{#1}} % to be removed
\newcommand{\green}[1]{\textcolor{green}{#1}} % to be discussed?
\newcommand{\yellow}[1]{\textcolor{yellow}{#1}} % ????
\newcommand{\magenta}[1]{\textcolor{magenta}{#1}} % warning - requires attention

% Define as cores dos links e informações do PDF
\makeatletter
\hypersetup{
    portuguese,
    colorlinks,
    linkcolor=black,
    citecolor=black,
    filecolor=black,
    urlcolor=black,
    breaklinks=true,
    pdftitle={\@title},
    pdfauthor={\@author},
    pdfsubject={\imprimirpreambulo},
    pdfkeywords={abnt, latex, abntex, abntex2}
}
\makeatother

% Redefinição de labels
\renewcommand{\algorithmautorefname}{Código}
\def\equationautorefname~#1\null{Equa\c c\~ao~(#1)\null}

% Cria o índice remissivo
\makeindex

% Início do documento
\begin{document}
    % Retira espaço extra obsoleto entre as frases.
    \frenchspacing
    % Elementos pré textuais
    \pretextual
    %
% Documento: Capa
%

\makeatletter
\begin{capa}
	\thispagestyle{empty}%limpa estilo da pagina
	\setlength{\baselineskip}{0.72\baselineskip}
    \begin{center} %Alinhamento centralizado
	

    \textbf{\expandafter\uppercase\expandafter{\imprimirinstituicao}}\\
    \textbf{\expandafter\uppercase\expandafter{\imprimirdepartamento}}\\
    \textbf{\expandafter\uppercase\expandafter{\imprimirprograma}}
    \textbf{\expandafter\uppercase\expandafter{\imprimirmodalidade}}\\
	\vspace*{5cm}%Espaçamento entre linhas
	\large\textbf{\expandafter\uppercase\expandafter{\imprimirtitulotb}}\\
	\vspace*{6cm}%Espaçamento entre linhas	
	\small\textbf{\expandafter\uppercase\expandafter{\imprimirnomeautor}}\\
	\vspace*{8.5cm}%%Espaçamento entre linhas
	\large\textbf{\expandafter\expandafter{\imprimirlocalapresentacao \\\imprimirdataapresentacao}}
	%\small\textbf{\expandafter\uppercase\expandafter{\imprimirdata}}\\
		
		
		
	\end{center} %Alinhamento centralizado
\end{capa}
\makeatother

	

              % Capa
    %
% Documento: FOLHA DE ROSTO
%

\makeatletter
\begin{folhaderosto}
	\thispagestyle{empty}%limpa estilo da pagina
	
    \begin{center}
    
		\small\textbf{\expandafter\uppercase\expandafter{\imprimirnomeautor}}\\
		\vspace*{6.2 cm}%Espaço entre linhas
		\normalsize\textbf{\expandafter\uppercase\expandafter{\imprimirtitulotb}}\\
		
    \end{center}
	
	\vspace*{2.35 cm}%Espaçamento entre linhas
		    \large%tamanho da fonte 
    		\hfill%Estica horizontamente  com espaços
	    	\begin{minipage}{8 cm}%Minipagina
	    		\begin{small} %Muda tamanho da fonte
	    		\setlength{\baselineskip}{0.7\baselineskip}
				
		    	{Trabalho de Conclusão de Curso apresentado junto ao {\imprimirprograma }
		    	{\imprimirmodalidade} do {\imprimirinstituicao}{ - }{\imprimirdepartamento},
		    	como requisito à obtenção do título de
		    	{\imprimirgrau }.}\\{
		    	}\\Orientador\\ \\
		    	{\imprimirtitulacaoorientador }{ }{\imprimirorientador.}\\\\ 
		    	%se houver coorientador, retirar o comentário da linha abaixo 
		    	%{\imprimirtitulacaocoorientador }{ }{\imprimircoorientador.} - Coorientador
				
				
				\end{small} %Muda tamanho da fonte
		    \end{minipage}%%Minipagina
		    	
		    \vspace*{6 cm}%Espaçamento entre linhas
		    
		    \begin{center} %Alinhamento centralizado
		    	\normalsize %Muda tamanho da fonte
	    		\textbf{\imprimirlocalapresentacao 
	    		\\\imprimirdataapresentacao}
	    	\end{center}%Alinhamento centralizado

\end{folhaderosto}
\makeatother
        % Folha de rosto
    %
% Documento: FOLHA APROVAÇÃO
%

\makeatletter
\begin{folhadeaprovacao}

\thispagestyle{empty}%limpa estilo da pagina
	
	\begin{center}
    
		\small\textbf{\expandafter\uppercase\expandafter{\imprimirnomeautor}}\\
		\vspace*{2.0 cm}%Espaço entre linhas
		\normalsize\textbf{\expandafter\uppercase\expandafter{\imprimirtitulotb}}
		
    \end{center}
	
	\vspace*{0.35 cm}%Espaçamento entre linhas
		    \large%tamanho da fonte 
    		\hfill%Estica horizontamente  com espaços
	    	\begin{minipage}{8 cm}%Minipagina
	    		\begin{small} %Muda tamanho da fonte
	    		\setlength{\baselineskip}{0.7\baselineskip}
				
		    	{Trabalho de Conclusão de Curso apresentado junto ao {\imprimirprograma }
		    	{\imprimirmodalidade} do {\imprimirinstituicao}{ - }{\imprimirdepartamento},
		    	como requisito à obtenção do título de
		    	{\imprimirgrau }.}\\{
		    	}\vspace*{0.6 cm}\\Orientador\\ \\
		    	{\imprimirtitulacaoorientador }{ }{\imprimirorientador.}\\\\
		    	%se houver coorientador, remover o comentário abaixo
		    	%{\imprimirtitulacaocoorientador }{ }{\imprimircoorientador.} - Coorientador			
				
				\end{small} %Muda tamanho da fonte
		    \end{minipage}%%Minipagina
		    	
		    \vspace*{0.6 cm}%Espaçamento entre linhas
		    
		    \large%%tamanho da fonte 
    		\hfill%%Estica horizontamente  com espaços
	    	 
		    
		    \normalsize %Muda tamanho da fonte
		    \vspace*{0.5 cm}%Espaçamento entre linhas
		    
			\begin{flushleft}%Alinhamento centralizado
			    {Aprovada em: \textbf{{\imprimirmesapresentacao}.}}\\
		 		%\textbf{Banca Examinadora}\\ %Negrito
						
				\vspace*{2 cm}%Espaçamento entre linhas
				\centering \rule{12 cm}{.1 mm}\\
				\centering {\imprimirtitulacaoorientador}{ }{\imprimirorientador} - Orientador\\
				
				%Se houver coorientador, remover os comentários
				%\vspace*{1 cm}%Espaçamento entre linhas
				%\rule{12 cm}{.1 mm}\\
				%{\imprimirtitulacaocoorientador}{ }{\imprimircoorientador} - Coorientador\\
			
				\vspace*{1 cm}%Espaçamento entre linhas
				\rule{12 cm}{.1 mm}\\
				\imprimirtitulacaoexamum\imprimirnmexamum{} - Avaliador \\
				\imprimirinstexamum  - \imprimirdepartamentoexamum\\
				
				\vspace*{1 cm}%Espaçamento entre linhas
				\rule{12 cm}{.1 mm}\\
				\imprimirtitulacaoexamdois\imprimirnomeexamdois{} - Avaliador \\
				\imprimirinstexamdois  - \imprimirdepartamentoexamdois
	
				\vspace*{1.3 cm}%Espaçamento entre linhas
		    \end{flushleft}%Alinhamento centralizado


\end{folhadeaprovacao}
\makeatother
    % Folha de aprovação
    %
% Documento: Dedicatória
%

\thispagestyle{empty}
\begin{dedicatoria}

Aqui você pode inserir uma homenagem ou dedicar seu trabalho.

\end{dedicatoria}       % Dedicatória
    %
% Documento: Agradecimento
%

\begin{agradecimento}
	
		\noindent Neste trecho você faz agradecimentos dirigidos aqueles que contribuíram de maneira relevante a elaboração do trabalho. Elemento opcional segundo a norma da ABNT NBR 14724 de 2011.

\end{agradecimento}

    % Agradecimentos
    %
% Documento: Epígrafe
%
\thispagestyle{empty}
\begin{epigrafe}
\begin{em}

''Aqui coloca-se o texto da epígrafe, que é um elemento opcional. Trata-se de uma citação curta que tenha relação com o tema estudado.''
\newline
\newline
\end{em}
\begin{autorepigrafe}
Fonte da epígrafe
\end{autorepigrafe}

\end{epigrafe}


          % Epígrafe
    %
% Documento: Resumo (Português)
%

\begin{RESUMO}
\thispagestyle{empty}
	
		\noindent O resumo consiste em um texto conciso e objetivo, de 150 a 500 palavras, que tem como objetivo apresentar o trabalho desenvolvido no TCC. Após o resumo, o aluno deve descrever um conjunto de palavras-chave que caracterizam o trabalho desenvolvido. Atualmente, a NBR 6028 (Associação Brasileira de Normas técnicas, 2003b) descreve esse elemento.

		\vspace*{0.5cm}\noindent\textbf{Palavras-chave}: Conclusão. Trabalho. Curso. NBR. ABNT.
		

\end{RESUMO}


          % Resumo na língua vernácula
    %
% Documento: Resumo (Inglês)
%

\begin{ABSTRACT}
	
		\noindent O resumo em língua estrangeira trata-se de uma tradução do resumo para uma língua estrangeira. No TCC do curso de Tecnologia em Análise e Desenvolvimento de Sistemas, a língua estrangeira deve ser obrigatoriamente o inglês. Nesse caso, esse elemento terá o título Abstract. Após o resumo em inglês, o aluno também deve descrever um conjunto de Keywords, que correspondem à tradução para a língua inglesa das palavras-chave especificadas no resumo.


		\vspace*{0.5cm}\noindent\textbf{Keywords}: Work. Course. NBR. ABNT.
		
\end{ABSTRACT}
          % Resumo em língua estrangeira
    %
% Documento: Lista de figuras
%

\pdfbookmark[0]{\listfigurename}{lof}
\listoffigures*
\cleardoublepage

      % Lista de figuras
    %
% Documento: Lista de tabelas
%

\pdfbookmark[0]{\listtablename}{lot}
\listoftables*
\cleardoublepage
      % Lista de tabelas
    %
% Documento: Lista de quadros
%

\pdfbookmark[0]{\listofquadrosname}{loq}
\listofquadros*
\cleardoublepage
      % Lista de quadros
    %
% Documento: Lista de algoritmos
%

\newcommand{\algoritmoname}{Algoritmo}
\renewcommand{\listalgorithmcfname}{Lista de algoritmos}

\floatname{algocf}{\algoritmoname}
\newlistof{listofalgoritmos}{loa}{\listalgoritmoname}
\newlistentry{algocf}{loa}{0}

\counterwithout{algocf}{chapter}
\renewcommand{\cftalgocfname}{\algoritmoname\space}
\renewcommand*{\cftalgocfaftersnum}{\hfill--\hfill}

\pdfbookmark[0]{\listalgorithmcfname}{loa}
\listofalgorithms
\cleardoublepage
      %lista de algoritimos
    %
% Documento: Lista de abreviaturas e siglas
%

\begin{siglas}
	\setlength{\baselineskip}{0.7\baselineskip}
	
    \item[ADS] Análise e Desenvolvimento de Sistemas
    \item[IFPB] Instituto Federal de Educação, Ciência e Tecnologia da Paraíba
    \item[NBR] Norma Brasileira
    \item[TCC] Trabalho de Conclusão do Curso
    
\end{siglas}
       % Lista de abreviaturas e siglas
    
    \sumario
\begin{OnehalfSpacing}

    % Elementos textuais
    \textual
    %
% Documento: Introdução
%
\chapter{INTRODUÇÃO}\label{chap:introducao}

Apresentar brevemente o documento.

\section{Problemática}

Apresentar dados que descrevam o problema que sua solução pretende resolver, em que linha de pesquisa atua. Dar preferência por dados oficiais, outros trabalhos relacionados. 

\section{Descrição da solução}

Apresentar, de maneira sucinta, a solução proposta. Descrever o que faz, comparar com outras ferramentas semelhantes.

\section{Objetivos}

\subsection{Objetivo Geral}

Descrever em um parágrafo qual o objetivo do seu trabalho.


\subsection{Objetivos Específicos}

Descrever quais objetivos secundários devem ser atingidos para chegar no objetivo geral.

\begin{itemize}
    \item Objetivo 1
    \item Objetivo 2
    \item Objetivo 3
\end{itemize}

\section{Atividades}

Descrever quais serão as atividades a serem desenvolvidas no trabalho. Lembrar que não há somente a parte de desenvolvimento, as tarefas de pesquisar, escrever, entram aqui.

\section{Cronograma}

Apresentar o cronograma das atividades, especialmente para o TCC 2, pois é uma forma de mostrar a viabilidade. Pode organizar as atividades em meses.

\section{Organização do documento}

Apresentar quais são os capítulos seguintes do trabalho e o que é abordado em cada um deles.            % Introdução
    \chapter{Fundamentação teórica} \label{chap:fundamentacao}

Apresentar aqui todos os conceitos teóricos relacionados ao trabalho e que não sejam triviais.
           % Orientações
    \chapter{Trabalho}\label{chap:trabalho}

\section{Análise}

Apresentar os principais elementos de análise: definição de stakeholders, requisitos, casos de uso, protótipos. 


\section{Projeto}

A depender de cada projeto: arquitetura do sistema, projeto do banco de dados, diagramas UML.

\section{Tecnologias utilizadas}

Apresentar e referenciar as tecnologias que foram utilizadas. Não esquecer de especificar as versões.

\section{Resultados parciais}

Mostrar o que já está implementado.
          % Apresentação
    \chapter{Considerações finais}\label{chap:consideracoes}

Apresentar os resultados obtidos até o momento, ver quais objetivos foram atingidos, quais serão as próximas etapas.

\section{Exemplos de elementos}

\subsection{Exemplo de citação}

Citação implícita \cite{mittelbach2004latex} ou explícita \citeonline{mittelbach2004latex}.

Citação direta:

\begin{citacao}
Texto aqui, com mais de 3 linhas. \cite[pp. 10]{mittelbach2004latex}
\end{citacao}

\subsection{Exemplo de figura}

A Figura \ref{fig:exemplo} apresenta um modelo desse elemento.

\begin{center}
\begin{figure}[H]
	\vspace*{0,1cm}
	\centering
	\caption{Exemplo de figura}
    \fbox{\includegraphics[scale=1.5]{./04-figuras/campus-cajazeiras.png}}
    \label{fig:exemplo}
    {\fonte{Elaborado pelo autor}}\\
\end{figure}
\end{center}

\subsection{Exemplo de tabela}

Um exemplo de tabela é representado pela Tabela \ref{tab:exemplo}.

\begin{table}[h]
\centering
\caption{Exemplo de tabela}
\label{tab:exemplo}
%Essa tabela não representa nenhuma informação real, foi elaborada somente para servir de modelo
\vspace{0.5cm}
\begin{tabular}{l|lr}
 
Descrição & Valor \\ % Note a separação de col. e a quebra de linhas
\hline                               % para uma linha horizontal
Linha & 1 \\
Linha & 2 \\
Linha & 3 
\end{tabular}
\fonte{Elaborado pelo autor}
\end{table}

\subsection{Exemplo de quadro}

Um exemplo de quadro é apresentado no Quadro \ref{qua:exemplo}

\begin{quadro}[H]

	\begin{center}
	\caption{Exemplo de quadro
	\label{qua:exemplo}}
		\begin{tabular}{|p{5cm}|p{5cm}|}
			\hline
			\textbf{Coluna 1} & \textbf{Coluna 2}\\
			\hline
			Texto & Texto \\
			\hline
		\end{tabular}
	\end{center}
	\vspace*{-0,8cm}

	\fonte{Elaborado pelo autor}
	
\end{quadro}

\subsection{Exemplo de Algoritmo}

Um exemplo de código é apresentado no Algoritmo \ref{alg:hello}.

\begin{algocf}
  \caption{Hello World em java}
  \begin{lstlisting}
class HelloWorldApp {
  public static void main(String[] args) {
    //Display the string
    System.out.println("Hello World!");
  }
}
\end{lstlisting}
\label{alg:hello}
\fonte{Elaborado pelo autor}
\end{algocf}
    
\end{OnehalfSpacing}
    % Elementos pós textuais
    \postextual
    %
% Documento: Referências Bibliográficas
%

\bibliography{refbase}\label{chap:referencias}    % Geração automática das referências por meio do arquivo 'refbase.bib'
       % Referências
    \begin{OnehalfSpacing}
    %
% Documento: Apêndices
%

\begin{apendicesenv}

\chapter{Exemplo de apêndice}
\label{ape:exemplo}

Exemplo de texto de apêndice

\end{apendicesenv}
         % Apêndices
    %
% Documento: Anexos
%

\begin{anexosenv}

\chapter{Primeiro anexo}
\label{anexo:anexo1}
Exemplo do corpo do anexo

\end{anexosenv}            % Anexos
    \end{OnehalfSpacing}
    %\printindex                                             % Índice remissivo

\end{document}
